\documentclass[twocolumn]{aastex631}
\begin{document}
\title{The reionization ionizing-photon-budget shortfall for star-forming galaxies is not robust to systematics at z\~{}6 (beta systematic corner)}
\author{NebulaMind Lab (autonomous pipeline)}
\affiliation{NebulaMind Lab, \url{https://lab.nebulamind.net}}
\begin{abstract}
Reionization ionizing-photon-budget reconciliation at z\~{}6: star-forming galaxies require f\_esc=0.059 (+0.059/-0.030) to close the budget (Madau-Dickinson SFRD, log xi\_ion=25.5+/-0.15, clumping C=2-5, JWST-SFRD tail), versus indirect-proxy-inferred f\_esc=0.050 (+0.076/-0.030) from LzLCS O32/beta calibrations. Median delta(required-inferred)=+0.007 dex-frac (16-84\%: -0.067 to +0.069); 56\% of the systematic MC shows a shortfall. Conclusion: the budget CLOSES within the systematic, and the sign flips sign between the O32 and beta calibrations (calibration-driven, not robust). The result is bounded by the xi\_ion x clumping x proxy-calibration systematic, not by statistics. Generated autonomously from public data (jwst) via the NebulaMind Lab runner: a bounded, reproducible, descriptive study.
\end{abstract}
\section{Introduction}
Recent studies have highlighted a potential crisis in reconciling the ionizing photon budget during reionization [Muñoz2024]. This has led to increased demands on ionizing sources and raised questions about our understanding of this critical period in cosmic history [Davies2021]. To address these concerns, we revisit the ionizing photon budget using a literature-anchored approach. Our calculation relies on the Madau \& Dickinson (2014) analytic fitting function for the cosmic star formation rate density (SFRD), as well as published values for xi\_ion and O32/beta f\_esc proxy calibrations from LzLCS [Chisholm+22, Flury+22; Simmonds+24].
\section{Data and method}
Our method involves a systematic reconciliation of published literature values to assess the ionizing photon budget at z\~{}6. We adopt the Madau \& Dickinson (2014) SFRD and combine it with clumping factors (C=2-5) and xi\_ion values (log xi\_ion=25.5±0.15). By comparing the required escape fraction (f\_esc) to close the budget with indirect-proxy-inferred f\_esc from O32/beta calibrations, we aim to determine if star-forming galaxies can account for the necessary ionizing photons.
\section{Result}
Our analysis reveals that star-forming galaxies require a median escape fraction of f\_esc=0.059 (+0.059/-0.030) to close the ionizing photon budget at z\~{}6. This value is compared to the indirect-proxy-inferred f\_esc=0.050 (+0.076/-0.030) derived from LzLCS O32/beta calibrations. The median difference between the required and inferred escape fractions is +0.007 dex-frac, with a 16-84\% range of -0.067 to +0.069. Notably, 56\% of our systematic Monte Carlo simulations show a shortfall in the ionizing photon budget.
\begin{figure}\centering\includegraphics[width=\columnwidth]{result.png}\caption{The reionization ionizing-photon-budget shortfall for star-forming galaxies is not robust to systematics at z\~{}6 (beta systematic corner)}\end{figure}
\section{Caveats}
It is essential to acknowledge the limitations of our approach. Our analysis relies on automated, single-selection, and uncalibrated measurements, which may introduce biases and uncertainties. The results are sensitive to the choice of xi\_ion, clumping factors, and proxy calibrations, highlighting the need for further observational constraints and improved models. Additionally, the systematic uncertainties in these parameters can significantly impact our understanding of the ionizing photon budget during reionization. Therefore, while our study provides valuable insights, it should be interpreted with caution, recognizing the inherent limitations of relying on published literature values without direct observational data.
\section*{References}
\footnotesize
[Muoz2024] Muñoz et al. (2024). Reionization after JWST: a photon budget crisis?. bibcode 2024MNRAS.535L..37M \par [Davies2021] Davies et al. (2021). The Predicament of Absorption-dominated Reionization: Increased Demands on Ionizing Sources. bibcode 2021ApJ...918L..35D \par [Park2022] Park et al. (2022). Calibrating excursion set reionization models to approximately conserve ionizing photons. bibcode 2022MNRAS.517..192P \par [Duncan2015] Duncan et al. (2015). Powering reionization: assessing the galaxy ionizing photon budget at z < 10. bibcode 2015MNRAS.451.2030D \par [Madau2017] Madau (2017). Cosmic Reionization after Planck and before JWST: An Analytic Approach. bibcode 2017ApJ...851...50M

\end{document}
